\chapter{Introduzione}
\label{cap:introduzione}

Il sito web e la documentazione di un ateneo raccolgono un'enorme quantità di informazioni --- regolamenti, tasse,
scadenze, piani di studio, materiali didattici --- ma lo fanno in modo estremamente dispersivo: argomenti correlati
tra loro finiscono su pagine diverse, i documenti sono distribuiti su sezioni distinte del sito senza un percorso di
consultazione unico, e le informazioni cambiano da un anno accademico all'altro senza che sia sempre immediato
capire dove si trovi la versione aggiornata. Uno studente che cerca una risposta puntuale (\emph{``quando scade la
prima rata?''}) o composita (\emph{``quali sono le diverse modalità di pagamento?''}) si trova spesso a dover
incrociare più fonti per ricostruire un quadro coerente, con un dispendio di tempo non indifferente e un margine
concreto di errore di interpretazione.

A questa dispersione si aggiunge una seconda difficoltà, che riguarda gli studenti che per ragioni diverse faticano a
seguire le lezioni con regolarità: per loro, un assistente capace di rispondere anche su contenuti didattici --- non
solo su procedure amministrative --- rappresenta un supporto concreto per recuperare e comprendere quanto trattato a
lezione, anche senza avervi potuto assistere di persona. Lo stesso tipo di supporto porta con sé un beneficio
speculare per i docenti: chi tiene più corsi, spesso con un numero elevato di studenti iscritti, fatica a propria
volta a rispondere individualmente a tutte le domande che riceve, molte delle quali ricorrenti e già risposte in
precedenza ad altri studenti. Un assistente che risponde autonomamente alle domande più comuni allegerisce questo
carico, lasciando ai docenti le domande che richiedono davvero la loro attenzione diretta.

La soluzione proposta in questo lavoro di tesi è un assistente conversazionale basato sul paradigma del
\emph{Retrieval-Augmented Generation} (RAG). A fronte di una domanda posta in linguaggio naturale, il sistema
recupera automaticamente, tra tutta la documentazione e il materiale didattico indicizzato, i frammenti più
pertinenti alla domanda posta, e li fornisce a un modello linguistico incaricato di formulare una risposta diretta
basandosi su di essi, citando la fonte da cui ogni affermazione proviene. In questo modo chi pone la domanda --- sia
esso uno studente o un docente che voglia verificare rapidamente un'informazione --- non deve più cercare manualmente
tra le pagine del sito o i materiali del corso: riceve una risposta puntuale e, soprattutto, verificabile, perché
ancorata a un documento reale e non semplicemente generata liberamente dal modello.

Il presente lavoro ha natura di \textbf{studio di fattibilità}: l'obiettivo primario non è consegnare un sistema
pronto per il deployment in produzione, ma verificare se un assistente di questo tipo --- completamente locale,
vincolato a un hardware non professionale --- possa rispondere con una qualità sufficiente a domande reali in questo
dominio, e individuare concretamente dove e perché un simile approccio fatica o fallisce. I contributi di questo
lavoro sono, di conseguenza: la progettazione e l'implementazione di un assistente RAG completamente locale per il
dominio universitario; un dataset diagnostico costruito appositamente per isolare le cause di errore lungo l'intera
pipeline, dal recupero alla generazione; un confronto sperimentale tra più modelli linguistici open-weight eseguibili
sull'hardware disponibile; l'identificazione e la risoluzione di una serie di problematiche concrete emerse durante
lo sviluppo, documentate insieme al processo diagnostico che le ha portate alla luce.

Il resto di questo lavoro di tesi è organizzato come segue. Il Capitolo~\ref{cap:stato-arte} inquadra il problema
nella letteratura esistente, introduce i concetti alla base dei modelli linguistici di grandi dimensioni e del
paradigma RAG, e passa in rassegna sistemi esistenti che affrontano problemi analoghi in ambito universitario. Il
Capitolo~\ref{cap:architettura} discute i vincoli di progetto che hanno guidato le scelte realizzative e descrive in
dettaglio l'architettura del sistema sviluppato, le scelte tecnologiche operate per ciascun componente e le
problematiche concrete incontrate durante lo sviluppo, insieme alle soluzioni adottate per risolverle. Il
Capitolo~\ref{cap:validazione} presenta la metodologia di valutazione sperimentale impiegata, gli esperimenti
condotti e i risultati ottenuti, inclusi alcuni casi di errore analizzati in dettaglio. Il
Capitolo~\ref{cap:conclusioni}, infine, riassume il contributo di questo lavoro rispetto agli obiettivi qui
enunciati, ne discute i limiti noti e ne delinea i possibili sviluppi futuri.
